\documentclass[10pt,a4paper]{amsart}
\usepackage[margin=19mm]{geometry}
\usepackage{amsmath,amssymb,mathtools}
\usepackage[scr=rsfs,cal=euler]{mathalfa}
\usepackage{xfrac}
\usepackage[unicode,hidelinks,bookmarks=false]{hyperref}
\newcommand{\ie}{i.e.\ }
\newcommand{\cf}{cf.\ }
\newcommand{\resp}{respectively\ }
\newcommand{\Subsection}{\subsection}
\providecommand{\nobreakdash}{\nobreak\hbox{-}}
\setlength{\emergencystretch}{2em}
\multlinegap0pt
\usepackage{bm}
\usepackage{bbm}
\usepackage{dsfont}
\usepackage{xspace}
\usepackage{cancel}
\usepackage{mathtools}
\usepackage{amsthm}
\makeatletter
\@ifundefined{theorem}{\newtheorem{theorem}{Theorem}}{}
\@ifundefined{lemma}{\newtheorem{lemma}[theorem]{Lemma}}{}
\makeatother
\DeclareMathOperator{\expo}{expo}
\newcommand{\RR}{\mathbb{R}}
\newcommand{\ZZ}{\mathbb{Z}}
\newcommand{\ball}{\mathbb{B}}
\newcommand{\origin}{o}
\title[Minimal covering bodies and Brunn-Minkowski type inequalitie]{Minimal covering bodies and Brunn-Minkowski type inequalities for the covering radius}
\author{Giulia Codenotti, Ansgar Freyer, Katarina Krivoku\'ca}
\date{Source: arXiv:2606.14442v1 (CC BY 4.0)}
\begin{document}
\maketitle

\noindent\emph{Source and licence.} Minimal covering bodies and Brunn-Minkowski type inequalities for the covering radius. Version arXiv:2606.14442v1 (CC BY 4.0), distributed under the Creative Commons Attribution 4.0 International licence, as recorded in the arXiv licence metadata for this exact version and shown on the abstract page https://arxiv.org/abs/2606.14442v1. Full rights notice: licenses/arxiv-2606.14442-CC-BY-4.0.txt.\par\smallskip

\noindent\textit{Editorial scope.} Counted unit: the Lemma whose source label is \texttt{lemma:balls\_lemma} in the article (source0.tex lines 407--423, source label \texttt{lemma:balls\_lemma}), delivered together with the article's own complete original proof of it (source0.tex lines 425--471, one \texttt{proof} environment of 47 lines). No mathematics is written, repaired or re-derived: statement and proof are the article's own text line for line. Required context retained uncounted: none. Every definition, display and label of those lines is retained unchanged. Omitted from this package and not redistributed: 1--406, 424--424, 472--1230. Nothing inside a retained excerpt is omitted or altered: every retained body is delivered exactly as the article's lines are written, so no omission receipt of any kind accompanies this package. Citations: the retained excerpts carry no citation command at all, so no bibliography entry of the article is transcribed and no reference is redistributed. Reference adaptation: none was needed and none was made; every reference inside the delivered unit resolves inside it, so no unresolved reference or citation is produced. Printed slips kept literal: no printed word, symbol, hypothesis or number of the counted statement or of its proof was altered. Layout and class: the article's own class is \texttt{amsart} with packages including xcolor, tikz, amsmath, amssymb, amsthm, amsfonts, graphicx, subcaption, cprotect, cite, fontenc, enumerate, todonotes, changes; the wrapper is the lane's pinned \texttt{amsart} editorial wrapper at 10pt on A4, which restates the theorem-like environments the retained text needs and adds the article's own macro definitions that the retained text needs (\texttt{\textbackslash RR}, \texttt{\textbackslash ZZ}, \texttt{\textbackslash ball}, \texttt{\textbackslash expo}, \texttt{\textbackslash origin}), with extra wrapper packages (bm, bbm, dsfont, xspace, cancel, mathtools). The delivered numbering is the unit's own. Assets: no retained span contains a figure environment, a graphics-inclusion command, a PDF-page inclusion, a table or a TikZ picture; no image file is copied into this package, the package declares no asset dependency and no image file is redistributed with it. Licence: version arXiv:2606.14442v1 of this article is distributed under the Creative Commons Attribution 4.0 International licence, as recorded in the arXiv licence metadata for this exact version and shown on the abstract page https://arxiv.org/abs/2606.14442v1. A full rights notice is delivered at licenses/arxiv-2606.14442-CC-BY-4.0.txt. Characters: every retained excerpt is pure ASCII, so the delivered engine, XeLaTeX with its default Latin Modern text font, reports no missing character.
\begin{lemma}
\label{lemma:balls_lemma} Let $K\subset \RR^d$ be a convex body. Then, the following statements are equivalent:
\begin{enumerate}
\item $K$ is minimal covering.
\item $K$ is covering and for all $v\in \expo(K)$, $u\in \RR^d$ exposing $v$ in $K$ and all $\varepsilon>0$, the convex body
$$K_{u, \varepsilon}:=K\cap \{x\in \RR^d \; | \; \langle x,u\rangle \leq \langle v,u\rangle-\varepsilon \}$$
is not covering.
\item For all $v\in \RR^d$ and $0<\varepsilon<\frac{1}{2}$, 
\begin{equation}\label{eq:balls_covering}
\bigcup \limits_{z\in \ZZ^d} (K-z) \cap \ball (v, \varepsilon)=\ball(v,\varepsilon),
\end{equation}
and, if $v\in\expo(K)$,
\begin{equation}\label{eq:balls_condition}
\ball(v,\varepsilon)\cap K \not\subseteq\bigcup \limits_{z\in \ZZ^d\setminus \{\origin\}} (K-z) \cap \ball (v, \varepsilon).
\end{equation}
\end{enumerate}
\end{lemma}

\begin{proof} 
We will first prove the equivalence of the first two statements.
For an arbitrary $v\in \expo(K)$, $u\in \RR^d$ exposing $v$ in $K$ and all $\varepsilon>0$, notice that $K_{u,\varepsilon}\subsetneq K$. Therefore, if $K$ is minimal covering, this convex body will not be covering.

For the converse implication, assume that $K$ is not minimal covering. Then there exists $L\subsetneq K$ such that $L+\ZZ^d=\RR^d$. Let $v$ be an exposed point in $K$ which is not in $L$ and $u\in \RR^d$ a functional exposing $v$ in $K$. For $\varepsilon:=\langle v,u\rangle-\max_{x\in L}\langle x,u\rangle>0$, notice that $L\subset K_{u,\varepsilon}$. Therefore, $\RR^d=L+\ZZ^d\subset K_{u,\varepsilon}+\ZZ^d,$ and the second statement cannot hold.

The next goal is to show that the second statement implies the third. 
Let $v\in \RR^d$ and $0<\varepsilon<\frac{1}{2}$ be arbitrary. Since $K+\ZZ^d=\RR^d$, we can intersect both sides with $\ball(v, \varepsilon)$ and writing out the Minkowski sum we obtain 
\[
\bigcup \limits_{z\in \ZZ^d} \ball (v, \varepsilon)\cap (K-z)=\ball(v,\varepsilon).
\]
 Towards a contradiction, assume that there is an exposed point $v\in \expo(K)$ and $0<\varepsilon<\frac{1}{2}$ such that
\begin{equation}\label{eq:balls_contradiction}
\ball(v,\varepsilon)\cap K \subseteq\bigcup \limits_{z\in \ZZ^d\setminus \{\origin\}} \ball (v, \varepsilon)\cap (K-z).
\end{equation}
\noindent Let $u\in \RR^d$ be a vector exposing $v$ in $K$, and set 
\[
\varepsilon':= \langle v,u\rangle-\max_{x\in K\cap \partial \ball(v,\varepsilon)}\langle x,u\rangle>0.
\]
%We claim that then $K_{u, \varepsilon'}$ will be covering.\\
Let $y\in K\setminus K_{u, \varepsilon'}$. Then, \[
\langle y,u\rangle \geq \langle v,u\rangle-\varepsilon'=\max_{x\in K\cap \partial \ball(v,\varepsilon)}\langle x,u\rangle.
\]
 Let $y'$ be the intersection of $\partial\ball(v, \varepsilon)$ and the ray starting from $v$ going through $y$. Then, $y=v+(y-v)=v+\lambda(y'-v)$, where $\lambda>0$ by definition and $\lambda\leq 1$ if and only if $y\in \ball(v, \varepsilon)$.
Calculating
\[
\begin{split}
\langle y,u\rangle &= \langle v+\lambda(y'-v) , u \rangle = (1-\lambda)\langle v,u\rangle +\lambda \langle y',u\rangle\\
& \leq (1-\lambda)\langle v,u\rangle + \lambda\max_{x\in K\cap \partial \ball(v,\varepsilon)} \langle x,u\rangle \\
&=(1-\lambda)\langle v,u\rangle+\lambda(\langle v,u\rangle-\varepsilon')=\langle v,u\rangle-\lambda\varepsilon'.
\end{split}
\]
Since $\langle y,u\rangle\geq\langle v,u\rangle - \varepsilon'$ by supposition and $\varepsilon'>0$, we conclude that $\lambda\leq 1$, that is $y\in \ball(v,\varepsilon)$. Therefore, $K\setminus K_{u,\varepsilon'}\subseteq \ball(v,\varepsilon)\cap K$.

 Let $z\in \ZZ^d\setminus \{ \origin\}$. Notice that $(K\setminus K_{u, \varepsilon'})\cap \ball(v+z, \varepsilon)\subseteq \ball(v,\varepsilon)\cap \ball(v+z,\varepsilon)$. However, we chose $\varepsilon < \frac{1}{2}$ which provides $\ball(v, \varepsilon)\cap\ball(v+z,\varepsilon)=\emptyset$. Since $K_{u,\varepsilon'}\subset K$ we can conclude that $K\cap \ball(v+z, \varepsilon)\subseteq K_{u,\varepsilon'}$. Translating by $-z$, we get \(\ball(v,\varepsilon)\cap(K-z)\subseteq K_{u,\varepsilon'}-z\).
 
Combining these two inclusions with equation \eqref{eq:balls_contradiction}, we see
\[
\begin{split}
K\setminus K_{u,\varepsilon'}&\subseteq K\cap \ball(v,\varepsilon)\subseteq \bigcup\limits_{z\in \ZZ^d\setminus\{\origin\}}\ball(v,\varepsilon)\cap(K-z)\\
&\subseteq \bigcup\limits_{z\in \ZZ^d\setminus\{\origin\}}\ball(v,\varepsilon)\cap(K_{u,\varepsilon'}+z)=\ball(v,\varepsilon)\cap (K_{u,\varepsilon'}+(\ZZ^d\setminus\{\origin\})).
\end{split}
\]
Specifically, this shows that $K\setminus K_{u, \varepsilon'}\subseteq K_{u, \varepsilon'}+\ZZ^d$, and since $\origin\in \ZZ^d$, $K\subset K_{u, \varepsilon'}+\ZZ^d$. By supposition of $K$ being covering, this shows that $K_{u, \varepsilon'}$ is covering, which is a contradiction with the second statement. 

Finally, we show that the third statement implies the first one. It follows directly from \eqref{eq:balls_covering} that $K$ is covering. To see that $K$ is minimal covering, let $L\subsetneq K$ be a convex body and $v\in \expo(K)\setminus L$. Set $\varepsilon:=\min \{ \frac{1}{4}, \frac{1}{2}d(v,L)\}>0$, where $d(v,L)$ denotes the Euclidean distance of $v$ to $L$. From the condition \eqref{eq:balls_condition} we see that there exists some $y\in \ball(v,\varepsilon)\cap K$ which is not in $\ball(v,\varepsilon)\cap (K-z)$ for any $z\in \ZZ^d\setminus\{\origin\}$. Since $L\subseteq K$, this implies that $y\notin L-z$ for all $z\in \ZZ^d\setminus \{\origin\}$. Because $\varepsilon\leq\frac{1}{2}d(v,L)$, $K\cap \ball(v,\varepsilon)\subseteq K\setminus L$. Since $y\in K\cap \ball(v,\varepsilon)$, this shows that $y\notin L$. Therefore, $L$ is not covering.
\end{proof}

\end{document}
